\scanpage{212}
\mspage{4}
% archive sheet number: 211

\unsure{Donc|Pour} moyennant G0 à G4, il faut prouver ceci : si $X\in\mathrm{Ob}\,C$, alors il \deleted{existe famille de descente effective} \deleted{universelle}\inserted{les} $S_i\longrightarrow e$, \unsure{tels|telles} que les $X_{S_i}/S_i$ soient constants finis, recouvrent $e$ pour la topologie de la descente effective. Utilisant G2 et G4, on est

\deleted{Soit $e_i$ \unsure{distincts} de $e$, \deleted{\unsure{donné}}\inserted{sommand} direct de $e$} \deleted{[G3] \deleted{et G2, entrainent que un \unsure{isomorphisme} d'un} prenons les parties finies $J$ de $I$, on peut} \deleted{\deleted{\unsure{monomorphisme direct}} considérer \deleted{$\unsure{\pi_J}$} $\operatorname{Sup}_{i\in J}e_i$ (sommes sont disjointes} \deleted{et universelles, et produits fibrés existent), ils} \deleted{forment une famille croissante d'objets, donc stationnaire.} \deleted{Soit $e'$ la valeur stationnaire, et $e''$ son ``complémentaire''.} \deleted{Je dis que $e''=\emptyset$ i.e. $e'=e$.} \deleted{En effet, considérons $X_{e''}$ sur $e''$, alors}

ramené à prouver que si $X\longrightarrow Y$ ($Y$ sommand direct de $e$, mais peu importe) est tel que $Y$ connexe $\ne\emptyset$, \unsure{et} $X$ connexe, alors $\exists\,Y'\longrightarrow Y$ de descente eff. universelle, avec $X_{Y'}\simeq[1,n]_{Y'}$ ($n$ convenable). On \deleted{\unsure{regarde la démarche|en baissant les dimensions}}\inserted{procède par récurrence sur l'entier} \deleted{\unsure{N}} $n$ de G5 (majorant le rang], en reprenant la démonstration plus haut

\uline{Remarque} : il se pourrait que G5 soit surabondant.

Noter qu'en présence des autres axiomes, on peut le remplacer par :

G5$'$ \deleted{\unsure{Il ex}} Pour tt composante connexe $e_i$ de $e$, $\exists$ foncteur exact $F:C\longrightarrow(\mathrm{Ens}\,f]$ tel que $F(e_i)\unsure{\ne|=}\emptyset$

